\documentclass[11pt]{article}

% ── Paquetes ────────────────────────────────────────────────────────────────
\usepackage[utf8]{inputenc}
\usepackage[T1]{fontenc}
\usepackage[spanish,es-nodecimaldot]{babel}
\usepackage{geometry}
\usepackage{fancyhdr}
\usepackage{graphicx}
\usepackage{xcolor}
\usepackage{fontawesome5}
\usepackage{listings}
\usepackage{float}
\usepackage[most]{tcolorbox}
\usepackage{hyperref}
\usepackage{amsmath}
\usepackage{amssymb}
\usepackage{booktabs}
\usepackage{multirow}
\usepackage{array}
\usepackage{tabularx}
\usepackage{caption}
\usepackage{subcaption}
\usepackage{url}
\usepackage{parskip}
\usepackage{tikz}
\usepackage{pgfplots}
\usepackage{enumitem}
\pgfplotsset{compat=1.18}
\usetikzlibrary{positioning,arrows.meta,shapes.geometric,shapes.misc,shapes.multipart,fit,backgrounds,calc,shadows,decorations.pathreplacing}

% ── Rutas de imágenes ───────────────────────────────────────────────────────
\graphicspath{{./}{images/}}

% ── Colores ─────────────────────────────────────────────────────────────────
\definecolor{htmlblue}{RGB}{41, 128, 185}
\definecolor{htmlgreen}{RGB}{39, 174, 96}
\definecolor{codebg}{RGB}{248, 249, 250}
\definecolor{grayborder}{RGB}{220, 220, 220}
\definecolor{qcolor}{RGB}{56, 142, 60}
\definecolor{ccolor}{RGB}{25, 118, 210}
\definecolor{gblue}{RGB}{26, 115, 232}
\definecolor{gred}{RGB}{217, 48, 37}
\definecolor{gyellow}{RGB}{249, 171, 0}
\definecolor{ggreen}{RGB}{30, 142, 62}
\definecolor{awsorange}{RGB}{255, 153, 0}
\definecolor{inkgray}{RGB}{60, 64, 67}
\definecolor{softgray}{RGB}{241, 243, 244}
\definecolor{mapbg}{RGB}{232, 240, 233}
\definecolor{purple}{RGB}{142, 36, 170}

% ── Listings ────────────────────────────────────────────────────────────────
\lstset{
    language=Java,
    basicstyle=\ttfamily\footnotesize,
    backgroundcolor=\color{codebg},
    keywordstyle=\color{blue}\bfseries,
    stringstyle=\color{orange},
    commentstyle=\color{gray}\itshape,
    showstringspaces=false,
    numbers=left,
    numberstyle=\tiny\color{gray},
    numbersep=8pt,
    frame=lines,
    framesep=5pt,
    rulecolor=\color{grayborder},
    breaklines=true,
    breakatwhitespace=true,
    tabsize=4,
    literate={á}{{\'a}}1 {é}{{\'e}}1 {í}{{\'i}}1 {ó}{{\'o}}1 {ú}{{\'u}}1 {ñ}{{\~n}}1 {²}{{$^2$}}1
}

% ── Entornos personalizados ─────────────────────────────────────────────────
\newcounter{countActividad}
\newenvironment{actividad}[1]{%
    \stepcounter{countActividad}
    \vspace{1.5em}
    \noindent{\large\color{htmlblue}\faCode\ \textbf{Actividad \ifnum\value{countActividad}<10 0\fi\thecountActividad: #1}}
    \par\medskip
}{%
    \vspace{1em}\hrule\vspace{1em}
}

% Tarjeta de indicador (KPI)
\newcommand{\kpi}[4]{%
  \begin{tcolorbox}[enhanced, colback=white, colframe=#1, boxrule=0pt, leftrule=3pt,
    arc=2pt, width=\linewidth, left=6pt, right=4pt, top=4pt, bottom=4pt,
    drop fuzzy shadow=black!25]
    {\color{#1}\Large #2}\hfill{\LARGE\bfseries\color{inkgray}#3}\\[1pt]
    {\footnotesize\color{gray}#4}
  \end{tcolorbox}}

% Cajas de nota / idea clave
\newtcolorbox{nota}[1][Nota]{enhanced, colback=htmlblue!5, colframe=htmlblue, boxrule=0.6pt, arc=3pt,
  fonttitle=\bfseries, coltitle=white, colbacktitle=htmlblue, title={\faInfoCircle\ #1}}
\newtcolorbox{clave}[1][Idea clave]{enhanced, colback=htmlgreen!6, colframe=htmlgreen, boxrule=0.6pt, arc=3pt,
  fonttitle=\bfseries, coltitle=white, colbacktitle=htmlgreen, title={\faLightbulb\ #1}}
\newtcolorbox{alerta}[1][Atención]{enhanced, colback=gred!5, colframe=gred, boxrule=0.6pt, arc=3pt,
  fonttitle=\bfseries, coltitle=white, colbacktitle=gred, title={\faExclamationTriangle\ #1}}

% Estilos TikZ comunes
\tikzset{
  box/.style={draw=#1, fill=#1!8, rounded corners=4pt, thick, align=center, font=\small,
    minimum height=1.0cm, inner sep=5pt, drop shadow={opacity=0.15}},
  box/.default=gblue,
  sbox/.style={draw=#1, fill=white, rounded corners=3pt, align=center, font=\footnotesize,
    minimum height=0.75cm, inner sep=4pt},
  sbox/.default=gray,
  group/.style={draw=#1, dashed, rounded corners=8pt, thick, inner sep=10pt},
  arr/.style={-{Stealth[length=2.6mm]}, thick, color=inkgray},
  darr/.style={{Stealth[length=2.6mm]}-{Stealth[length=2.6mm]}, thick, color=inkgray},
  lbl/.style={font=\scriptsize, fill=white, inner sep=1.5pt, text=inkgray, align=center},
  decision/.style={diamond, aspect=2.2, draw=gyellow!90!black, fill=gyellow!12, thick, align=center,
    font=\footnotesize, inner sep=1pt},
}

% Entidad de diagrama ER: {color}{nombre}{campos separados por \\}
\newcommand{\entidad}[3]{%
  \begin{tabular}{@{}l@{}}
    \colorbox{#1!18}{\makebox[3.3cm][l]{\textbf{#2}}}\\[1pt]
    \begin{tabular}{@{\hspace{4pt}}l@{\hspace{4pt}}}#3\end{tabular}
  \end{tabular}}

% Captura real si existe en images/, si no un esquema dibujado
\newcommand{\capturaO}[3]{\IfFileExists{images/#1}{\includegraphics[width=#2]{#1}}{#3}}

% ── Márgenes y encabezados ──────────────────────────────────────────────────
\geometry{a4paper, top=2.5cm, bottom=2.5cm, left=2.5cm, right=2.5cm}
\pagestyle{fancy}
\fancyhf{}
\fancyhead[L]{\small\textcolor{gray}{GeoTrack — Aplicativo de rastreo espacial}}
\fancyhead[R]{\small\textcolor{gray}{UNSAAC 2026}}
\fancyfoot[R]{\thepage}
\renewcommand{\headrulewidth}{0.4pt}
\setlength{\headheight}{14pt}

% ── Hipervínculos ───────────────────────────────────────────────────────────
\hypersetup{
    colorlinks=true,
    linkcolor=htmlblue,
    citecolor=htmlgreen,
    urlcolor=htmlblue,
}

\captionsetup{font=small, labelfont={bf,color=htmlblue}}
\setlist{itemsep=2pt, topsep=4pt}

% ════════════════════════════════════════════════════════════════════════════
\begin{document}

% ── Portada ─────────────────────────────────────────────────────────────────
\begin{titlepage}
  \begin{center}
    \LARGE\textbf{UNIVERSIDAD NACIONAL SAN ANTONIO ABAD DEL CUSCO}\\[0.4cm]
    \Large\textsc{Facultad de Ingeniería Eléctrica, Electrónica, Informática y Mecánica}\\[0.2cm]
    \large\textsc{Escuela Profesional de Ingeniería Informática y de Sistemas}
  \end{center}
  \vspace{0.8cm}
  \begin{center}
    \includegraphics[width=13cm]{logosU.png}
  \end{center}
  \vspace{0.8cm}
  \begin{center}
    {\huge \textbf{Informe proyecto}}\\[0.4cm]
    {\LARGE \textbf{Aplicativo de rastreo espacial}}
  \end{center}
  \noindent\rule{\linewidth}{0.6pt}

  \vspace{1cm}
  \noindent\textbf{Curso:} Sistemas Embebidos\\[0.2cm]
  \textbf{Docente:} Jose Mauro Pillco Quispe\\[0.2cm]
  \textbf{Estudiante:}  Mogrovejo Ccorimanya Luis Angel (220961)
  \vfill
  \begin{center}
    {\Large \textsc{Cusco — Perú}}\\[0.3cm]
    {\Large \textsc{2026}}
  \end{center}
\end{titlepage}

% ── Contenido ───────────────────────────────────────────────────────────────

% ── Resumen ejecutivo ───────────────────────────────────────────────────────
\section*{Resumen ejecutivo}
\addcontentsline{toc}{section}{Resumen ejecutivo}

\textbf{GeoTrack} es un sistema integral de geolocalización en tiempo real compuesto por tres piezas que
trabajan juntas: una \textbf{aplicación móvil Android} (Flutter) que actúa como rastreador, un
\textbf{backend en la nube de AWS} (Node.js, Socket.IO y PostgreSQL) que recibe, almacena y retransmite las
posiciones, y un \textbf{centro de monitoreo web} (React y Leaflet) que muestra en vivo todos los
dispositivos, su historial y la reproducción animada de sus recorridos.

El teléfono funciona como un \emph{sistema embebido}: un sensor (receptor GNSS), un proceso residente que
sobrevive al cierre de la aplicación y al bloqueo de pantalla, almacenamiento local para operar sin
conexión, y un canal de comunicación que sincroniza con el servidor cuando hay red. Sobre las lecturas
crudas del GPS se implementó un \textbf{filtro de precisión} (validación por exactitud y por velocidad
física, filtro de Kalman con predicción por velocidad y rumbo, y fijación de posición estando quieto) que
eliminó los ``saltos'' a otros lugares y redujo un 80\,\% el movimiento falso con el usuario detenido.

\vspace{0.3cm}
\noindent
\begin{minipage}[t]{0.32\linewidth}\kpi{gblue}{\faMobile*}{1 s}{intervalo mínimo de captura}\end{minipage}\hfill
\begin{minipage}[t]{0.32\linewidth}\kpi{ggreen}{\faMapMarker*}{50\,785}{posiciones registradas en campo}\end{minipage}\hfill
\begin{minipage}[t]{0.32\linewidth}\kpi{gred}{\faRoute}{21}{sesiones de rastreo reales}\end{minipage}

\noindent
\begin{minipage}[t]{0.32\linewidth}\kpi{awsorange}{\faCloud}{2}{servicios AWS (EC2 + RDS)}\end{minipage}\hfill
\begin{minipage}[t]{0.32\linewidth}\kpi{purple}{\faCode}{3\,900}{líneas de código propio}\end{minipage}\hfill
\begin{minipage}[t]{0.32\linewidth}\kpi{gyellow!90!black}{\faChartLine}{$-80$\,\%}{deriva falsa estando quieto}\end{minipage}

\newpage
\tableofcontents
\newpage

% ════════════════════════════════════════════════════════════════════════════
\section{Introducción}

\subsection{Problema}
Rastrear la posición de personas o vehículos de forma continua es una necesidad común en logística,
seguridad personal, deporte y supervisión de personal en campo. Las soluciones comerciales suelen ser
cerradas, costosas o dependientes de una conexión permanente. Además, un teléfono móvil impone
restricciones propias de un sistema embebido: batería limitada, un sistema operativo que \emph{mata}
procesos en segundo plano para ahorrar energía, conectividad intermitente y un sensor GNSS cuya precisión
varía entre 3 y más de 100 metros según el entorno.

\subsection{Objetivo general}
Desarrollar y desplegar un sistema de geolocalización en tiempo real, de extremo a extremo, compuesto por
una aplicación móvil rastreadora, un backend en la nube y una plataforma web de monitoreo, con una
interfaz moderna inspirada en Google Maps.

\subsection{Objetivos específicos}
\begin{enumerate}
  \item Capturar la ubicación de forma persistente en segundo plano (intervalo configurable de 1 a 10\,s),
        incluso con la aplicación cerrada o el teléfono bloqueado.
  \item Garantizar que ningún dato se pierda sin conexión mediante una arquitectura \emph{offline-first}.
  \item Calcular telemetría en el dispositivo: velocidad, distancia, tipo de actividad y calorías.
  \item Visualizar en vivo múltiples dispositivos, consultar su historial y reproducir sus recorridos.
  \item Desplegar la solución en AWS con infraestructura reproducible mediante scripts.
  \item Mejorar la calidad de la señal GNSS mediante filtrado y suavizado de las lecturas.
\end{enumerate}

\subsection{Alcance}
El sistema se organizó en cinco fases: (1) entorno local y despliegue en un dispositivo físico,
(2) infraestructura en AWS, (3) aplicación móvil, (4) plataforma web y (5) innovaciones adicionales.
La aplicación se probó en un Samsung Galaxy A26 real, en la ciudad del Cusco.

% ════════════════════════════════════════════════════════════════════════════
\section{Requerimientos}

\begin{table}[H]
\centering
\small
\begin{tabularx}{\linewidth}{>{\bfseries}l X >{\centering\arraybackslash}p{1.6cm}}
\toprule
\textbf{ID} & \textbf{Requerimiento funcional} & \textbf{Estado} \\
\midrule
RF-01 & Mapa a pantalla completa con botón flotante \emph{Iniciar/Detener} rastreo & \textcolor{ggreen}{\faCheckCircle} \\
RF-02 & Rastreo persistente en segundo plano (app cerrada o pantalla bloqueada) & \textcolor{ggreen}{\faCheckCircle} \\
RF-03 & Almacenamiento local sin conexión y envío de la cola al recuperar la red & \textcolor{ggreen}{\faCheckCircle} \\
RF-04 & Historial en el dispositivo: sesiones, ruta dibujada y nombres personalizados & \textcolor{ggreen}{\faCheckCircle} \\
RF-05 & Telemetría: velocidad, calorías y clasificación caminata / alta velocidad & \textcolor{ggreen}{\faCheckCircle} \\
RF-06 & Panel web en vivo con todos los dispositivos activos y color distintivo & \textcolor{ggreen}{\faCheckCircle} \\
RF-07 & Trazado de rutas en tiempo real & \textcolor{ggreen}{\faCheckCircle} \\
RF-08 & Consultas históricas por dispositivo, fecha y hora & \textcolor{ggreen}{\faCheckCircle} \\
RF-09 & Reproductor de rutas con velocidades x10, x20 y x30 (se añadió x1 y x60) & \textcolor{ggreen}{\faCheckCircle} \\
RF-10 & Al menos tres funcionalidades innovadoras (se implementaron seis) & \textcolor{ggreen}{\faCheckCircle} \\
\bottomrule
\end{tabularx}
\caption{Requerimientos funcionales y su cumplimiento.}
\end{table}

\begin{table}[H]
\centering
\small
\begin{tabularx}{\linewidth}{>{\bfseries}l X}
\toprule
\textbf{ID} & \textbf{Requerimiento no funcional} \\
\midrule
RNF-01 & Diseño moderno y minimalista inspirado en Google Maps (tarjetas flotantes, paleta neutra). \\
RNF-02 & Tolerancia a fallos de red: cero pérdida de datos e idempotencia en el servidor. \\
RNF-03 & Tiempo real: las posiciones aparecen en el panel en cuanto llegan al servidor (WebSockets). \\
RNF-04 & Seguridad: autenticación de dispositivos, base de datos privada y cifrada, mínimo privilegio en red. \\
RNF-05 & Despliegue reproducible: toda la infraestructura se crea, actualiza y destruye con scripts. \\
RNF-06 & Eficiencia energética: un solo flujo de GPS y envío por lotes. \\
\bottomrule
\end{tabularx}
\caption{Requerimientos no funcionales.}
\end{table}

% ════════════════════════════════════════════════════════════════════════════
\section{Arquitectura del sistema}

La Figura~\ref{fig:arquitectura} muestra los tres componentes y cómo se comunican. El teléfono envía las
posiciones por HTTP (REST) al servidor en AWS; el servidor las guarda en PostgreSQL y las retransmite al
instante a todos los navegadores conectados mediante WebSockets (Socket.IO).

\begin{figure}[!htbp]
\centering
\begin{tikzpicture}[node distance=0.55cm and 1.2cm]
  % Móvil
  \node[box=gblue, minimum width=3.4cm] (ui) {\faMobile*\ \textbf{Interfaz}\\{\scriptsize mapa, historial, ajustes}};
  \node[box=gblue, minimum width=3.4cm, below=of ui] (bg) {\faSatelliteDish\ \textbf{Servicio en 2.º plano}\\{\scriptsize GPS + filtro + telemetría}};
  \node[box=gblue, minimum width=3.4cm, below=of bg] (sql) {\faDatabase\ \textbf{SQLite local}\\{\scriptsize puntos + cola (outbox)}};
  \node[group=gblue, fit=(ui)(bg)(sql), label={[font=\small\bfseries, text=gblue]above:App móvil (Flutter)}] (mob) {};
  \draw[darr] (ui) -- (bg);
  \draw[darr] (bg) -- (sql);

  % AWS
  \node[box=awsorange, minimum width=3.6cm, right=2.2cm of ui] (ngx) {\faServer\ \textbf{Nginx}\\{\scriptsize proxy inverso :80}};
  \node[box=awsorange, minimum width=3.6cm, below=of ngx] (api) {\faCode\ \textbf{API Node.js}\\{\scriptsize Express + Socket.IO}};
  \node[box=awsorange, minimum width=3.6cm, below=of api] (rds) {\faDatabase\ \textbf{RDS PostgreSQL 16}\\{\scriptsize privada y cifrada}};
  \node[group=awsorange, fit=(ngx)(api)(rds), label={[font=\small\bfseries, text=awsorange!80!black]above:Nube AWS (us-east-1)}] (aws) {};
  \draw[arr] (ngx) -- (api);
  \draw[darr] (api) -- node[lbl, right]{SQL / TLS} (rds);

  % Web
  \node[box=ggreen, minimum width=3.2cm, right=2.2cm of ngx] (dash) {\faDesktop\ \textbf{Dashboard}\\{\scriptsize React + Leaflet}};
  \node[box=ggreen, minimum width=3.2cm, below=of dash] (live) {\faBroadcastTower\ \textbf{Panel en vivo}\\{\scriptsize rutas y alertas}};
  \node[box=ggreen, minimum width=3.2cm, below=of live] (hist) {\faHistory\ \textbf{Historial}\\{\scriptsize consultas y playback}};
  \node[group=ggreen, fit=(dash)(live)(hist), label={[font=\small\bfseries, text=ggreen]above:Centro de monitoreo}] (web) {};

  % Enlaces
  \draw[arr] (bg.east) -- node[lbl, below=2pt]{HTTP REST\\\texttt{x-api-key}} (api.west);
  \draw[darr] (ngx.east) -- node[lbl, above=2pt]{WebSocket} (dash.west);
  \draw[arr, dashed] (api.east) -- node[lbl, below=2pt]{eventos} (live.west);

  % Teselas
  \node[sbox=gray, below=1.1cm of rds, minimum width=4cm] (osm) {\faGlobe\ Teselas de mapa (OpenStreetMap)};
  \draw[arr, dotted] (osm.west) -| (sql.south);
  \draw[arr, dotted] (osm.east) -| (hist.south);
\end{tikzpicture}
\caption{Arquitectura general de GeoTrack.}
\label{fig:arquitectura}
\end{figure}

\subsection{Tecnologías utilizadas}

\begin{table}[H]
\centering
\small
\begin{tabularx}{\linewidth}{l l X}
\toprule
\textbf{Capa} & \textbf{Tecnología} & \textbf{Por qué se eligió} \\
\midrule
\multirow{5}{*}{\textcolor{gblue}{\faMobile*}\ Móvil}
  & Flutter 3 / Dart & Un solo código, compilación nativa (AOT) y alto rendimiento gráfico. \\
  & \texttt{flutter\_background\_service} & Servicio en primer plano de Android en un \emph{isolate} propio. \\
  & \texttt{geolocator} & Acceso al proveedor de ubicación fusionado de Android (GNSS + Wi-Fi + celdas). \\
  & \texttt{sqflite} & Base de datos SQLite para operar sin conexión. \\
  & \texttt{flutter\_map} & Mapa vectorial/teselas sin depender de claves de pago. \\
\midrule
\multirow{3}{*}{\textcolor{awsorange}{\faServer}\ Backend}
  & Node.js + Express & API REST ligera y asíncrona. \\
  & Socket.IO & WebSockets con reconexión automática para el tiempo real. \\
  & PostgreSQL 16 & Integridad, índices por tiempo y restricción \texttt{UNIQUE} para idempotencia. \\
\midrule
\multirow{2}{*}{\textcolor{ggreen}{\faDesktop}\ Web}
  & React 18 + Vite & Interfaz reactiva y compilación rápida. \\
  & Leaflet / React-Leaflet & Mapas interactivos con polilíneas, círculos y marcadores. \\
\midrule
\multirow{2}{*}{\textcolor{awsorange!80!black}{\faCloud}\ Nube}
  & AWS EC2 + RDS & Servidor con conexiones persistentes y base de datos gestionada con backups. \\
  & AWS CLI + PowerShell & Infraestructura reproducible por scripts. \\
\bottomrule
\end{tabularx}
\caption{Pila tecnológica.}
\end{table}

% ════════════════════════════════════════════════════════════════════════════
\section{Infraestructura en la nube (AWS)}

\begin{figure}[!htbp]
\centering
\begin{tikzpicture}[node distance=0.8cm and 1.3cm]
  \node[box=gblue, minimum width=2.6cm] (phone) {\faMobile*\ Teléfonos};
  \node[box=ggreen, below=0.5cm of phone, minimum width=2.6cm] (browser) {\faDesktop\ Navegadores};
  \node[box=purple, below=0.5cm of browser, minimum width=2.6cm] (admin) {\faLaptop\ PC administrador\\{\scriptsize scripts PowerShell}};
  \node[box=inkgray, right=1.0cm of browser, minimum width=2.0cm] (inet) {\faGlobe\\\textbf{Internet}};

  \node[box=awsorange, right=1.9cm of inet, yshift=1.3cm, minimum width=3.4cm] (eip) {\textbf{IP elástica}\\{\scriptsize dirección fija}};
  \node[box=awsorange, below=of eip, minimum width=3.4cm] (ec2) {\faServer\ \textbf{EC2 t3.small}\\{\scriptsize Amazon Linux 2023}\\{\scriptsize Nginx + Node.js (systemd)}};
  \node[box=awsorange, below=1.4cm of ec2, minimum width=3.4cm] (rds) {\faDatabase\ \textbf{RDS db.t4g.micro}\\{\scriptsize PostgreSQL 16, 20 GB}\\{\scriptsize cifrada, backups 7 días}};

  \node[group=gblue, fit=(ec2), inner sep=6pt, label={[font=\scriptsize, text=gblue]right:\begin{tabular}{l}\textbf{SG-API}\\TCP 80 \faArrowLeft\ todos\\TCP 22 \faArrowLeft\ IP admin\end{tabular}}] (sgapi) {};
  \node[group=gred, fit=(rds), inner sep=6pt, label={[font=\scriptsize, text=gred]right:\begin{tabular}{l}\textbf{SG-DB}\\TCP 5432 \faArrowLeft\ solo SG-API\end{tabular}}] (sgdb) {};
  \node[group=awsorange, fit=(eip)(sgapi)(sgdb), inner sep=14pt, label={[font=\small\bfseries, text=awsorange!80!black]above:VPC (us-east-1)}] (vpc) {};

  \draw[arr] (phone.east) -- (inet);
  \draw[darr] (browser) -- (inet);
  \draw[arr, purple, dashed] (admin.east) -- node[lbl, below right=0pt]{SSH :22} (inet);
  \draw[darr] (inet) -- node[lbl, above left=0pt]{HTTP :80} (eip.west);
  \draw[arr] (eip) -- (ec2);
  \draw[darr] (ec2) -- node[lbl, right]{TLS :5432} (rds);
\end{tikzpicture}
\caption{Despliegue en AWS con segmentación de red por Security Groups.}
\label{fig:aws}
\end{figure}

\begin{clave}[Decisión de diseño: EC2 en lugar de API Gateway + Lambda]
Socket.IO necesita \textbf{conexiones persistentes} con cada navegador. Un servidor EC2 detrás de Nginx las
mantiene de forma natural, con un costo predecible y una arquitectura simple de operar. La base de datos se
dejó en RDS para delegar en AWS los backups, los parches y el cifrado.
\end{clave}

\subsection{Seguridad de la red}
\begin{itemize}
  \item \textbf{Mínimo privilegio:} la base de datos no tiene IP pública; solo acepta conexiones desde el
        Security Group de la API.
  \item \textbf{SSH restringido} a la IP pública del administrador (regla \texttt{/32}).
  \item \textbf{Cifrado en reposo} del almacenamiento de RDS y conexión TLS entre la API y la base de datos.
  \item \textbf{Secretos fuera del repositorio:} token del panel, API key de dispositivos y URL de la base de
        datos se guardan en \texttt{infra/.deploy.json}, ignorado por git.
\end{itemize}

\subsection{Costo estimado}
\begin{table}[H]
\centering
\small
\begin{tabular}{l l r}
\toprule
\textbf{Recurso} & \textbf{Tipo} & \textbf{US\$/mes} \\
\midrule
EC2 & t3.small (2 vCPU, 2 GB) & $\approx 15$ \\
RDS PostgreSQL & db.t4g.micro + 20 GB gp2 & $\approx 13$ \\
Almacenamiento, IP elástica y transferencia & — & $\approx 2$ \\
\midrule
\textbf{Total} & & $\boldsymbol{\approx 30}$ \\
\bottomrule
\end{tabular}
\caption{Costo mensual aproximado (puede ser menor dentro de la capa gratuita).}
\end{table}

% ════════════════════════════════════════════════════════════════════════════
\section{Aplicación móvil: el rastreador}

\subsection{Organización del código}

La aplicación separa la interfaz del trabajo pesado. El rastreo vive en un \emph{isolate} independiente
(un hilo de Dart con memoria propia) dentro de un \textbf{servicio en primer plano} de Android, de modo que
cerrar la interfaz no detiene la captura.

\begin{figure}[!htbp]
\centering
\begin{tikzpicture}[node distance=0.35cm and 0.5cm]
  % UI
  \node[sbox=gblue, minimum width=2.6cm] (home) {\texttt{home\_screen}\\{\tiny mapa + FAB + SOS}};
  \node[sbox=gblue, minimum width=2.6cm, below=of home] (histo) {\texttt{history\_screen}\\{\tiny sesiones}};
  \node[sbox=gblue, minimum width=2.6cm, below=of histo] (det) {\texttt{session\_detail}\\{\tiny ruta + renombrar}};
  \node[sbox=gblue, minimum width=2.6cm, below=of det] (set) {\texttt{settings\_screen}\\{\tiny servidor, peso, intervalo}};
  \node[group=gblue, fit=(home)(set), label={[font=\small\bfseries, text=gblue]above:Isolate de interfaz}] (uig) {};

  % Background
  \node[sbox=ggreen, minimum width=3.0cm, right=2.4cm of home] (trk) {\texttt{tracking\_service}\\{\tiny GPS, watchdog, notificación}};
  \node[sbox=ggreen, minimum width=3.0cm, below=of trk] (flt) {\texttt{location\_filter}\\{\tiny precisión + Kalman}};
  \node[sbox=ggreen, minimum width=3.0cm, below=of flt] (tel) {\texttt{telemetry}\\{\tiny velocidad, actividad, kcal}};
  \node[sbox=ggreen, minimum width=3.0cm, below=of tel] (syn) {\texttt{sync\_service}\\{\tiny cola \faArrowRight\ servidor}};
  \node[group=ggreen, fit=(trk)(syn), label={[font=\small\bfseries, text=ggreen]above:Isolate en segundo plano}] (bgg) {};

  % Data
  \node[sbox=purple, minimum width=2.6cm, right=2.0cm of flt] (db) {\texttt{local\_db}\\{\tiny SQLite}};
  \node[sbox=purple, minimum width=2.6cm, below=of db] (cfg) {\texttt{config}\\{\tiny SharedPreferences}};
  \node[group=purple, fit=(db)(cfg), label={[font=\small\bfseries, text=purple]above:Datos}] (dg) {};

  \draw[darr] (uig.east |- trk) -- node[lbl, above]{\texttt{invoke('update')}\\\texttt{on('sos')}} (bgg.west |- trk);
  \draw[darr] (bgg.east |- db) -- (dg.west |- db);
  \draw[darr, dashed] (uig.south) |- ++(0,-0.6) -| node[lbl, pos=0.25]{lectura de historial y ajustes} (dg.south);
\end{tikzpicture}
\caption{Módulos de la app y comunicación entre isolates.}
\end{figure}

\subsection{Rastreo persistente en segundo plano}
Android cierra agresivamente los procesos en segundo plano, y más aún los fabricantes como Samsung. Para
evitarlo se combinaron varias técnicas:
\begin{itemize}
  \item \textbf{Servicio en primer plano} con \texttt{foregroundServiceType="location"} y notificación
        permanente que muestra velocidad, distancia, calorías y puntos en cola.
  \item \texttt{stopWithTask=false}: el servicio sigue vivo aunque se deslice la app fuera de recientes.
  \item \texttt{autoStartOnBoot}: si el teléfono se reinicia con una sesión activa, el rastreo se reanuda.
  \item Permiso de ubicación \emph{``Permitir todo el tiempo''} y exclusión de la optimización de batería.
  \item \textbf{Watchdog} cada 30\,s: si el flujo de GPS deja de entregar posiciones (algunos fabricantes lo
        pausan), fuerza una lectura y reinicia el flujo.
\end{itemize}

\subsection{Arquitectura \emph{offline-first} y sincronización}

Cada posición se guarda \textbf{primero en SQLite} y solo después se intenta enviar. Las operaciones que no
son puntos (iniciar/detener sesión, renombrar, SOS) van a una tabla \texttt{outbox} y se envían en orden.

\begin{figure}[!htbp]
\centering
\begin{tikzpicture}[x=1cm, y=-0.72cm, font=\footnotesize]
  \foreach \x/\n/\c/\i in {0/GPS/gblue/\faSatelliteDish, 3.2/App (2.º plano)/ggreen/\faCogs, 6.6/SQLite/purple/\faDatabase, 10.2/Servidor AWS/awsorange/\faServer, 13.4/Dashboard/ggreen/\faDesktop}{
    \node[box=\c, minimum width=2.4cm, minimum height=0.8cm] at (\x,0) {\i\ \textbf{\n}};
    \draw[gray!60, dashed] (\x,0.7) -- (\x,13.6);
  }
  \draw[arr] (0,1.4) -- node[lbl, above]{posición cada 1 s} (3.2,1.4);
  \draw[arr] (3.2,2.4) -- node[lbl, above]{\texttt{INSERT points (synced=0)}} (6.6,2.4);
  \node[lbl, fill=gred!10, draw=gred, rounded corners=2pt] at (8.4,3.6) {\faWifi\ sin red: el envío falla, los datos esperan};
  \draw[arr, gred] (3.2,4.8) -- node[lbl, above]{POST /api/locations} (9.7,4.8);
  \node[gred, font=\Large] at (9.9,4.8) {\faTimes};
  \draw[arr] (0,5.9) -- node[lbl, above]{posición} (3.2,5.9);
  \draw[arr] (3.2,6.6) -- node[lbl, above]{más puntos en cola} (6.6,6.6);
  \node[lbl, fill=ggreen!10, draw=ggreen, rounded corners=2pt] at (8.4,7.9) {\faWifi\ vuelve la red};
  \draw[arr] (6.6,9.0) -- node[lbl, above]{lote de hasta 500 puntos} (3.2,9.0);
  \draw[arr, ggreen] (3.2,10.0) -- node[lbl, above]{POST /api/locations} (10.2,10.0);
  \draw[arr] (10.2,10.0) -- ++(0.7,0) |- (10.2,10.8);
  \node[lbl, anchor=west, font=\tiny\ttfamily, align=left] at (11.0,10.4) {ON CONFLICT\\DO NOTHING};
  \draw[arr, dashed] (10.2,11.5) -- node[lbl, above]{\texttt{200 OK}} (3.2,11.5);
  \draw[arr] (3.2,12.3) -- node[lbl, above]{\texttt{synced = 1}} (6.6,12.3);
  \draw[arr] (10.2,12.8) -- node[lbl, above]{\texttt{emit('locations')}} (13.4,12.8);
\end{tikzpicture}
\caption{Diagrama de secuencia: corte de red y recuperación sin pérdida de datos.}
\label{fig:sync}
\end{figure}

\begin{clave}[Idempotencia]
La tabla \texttt{locations} tiene la restricción \texttt{UNIQUE(device\_id, ts)} y la inserción usa
\texttt{ON CONFLICT DO NOTHING}. Si un lote llega dos veces (por ejemplo, el servidor lo guardó pero la
respuesta se perdió por un corte de red), los puntos repetidos se ignoran: los reintentos nunca duplican
datos.
\end{clave}

\begin{lstlisting}[language=SQL, caption={Inserción idempotente de un lote de posiciones.}]
INSERT INTO locations (device_id, session_id, ts, lat, lng, speed, accuracy, ...)
VALUES ($1, $2, $3, ...), ($12, $13, $14, ...), ...
ON CONFLICT (device_id, ts) DO NOTHING
RETURNING *;
\end{lstlisting}

\subsection{Filtro de precisión de la señal GNSS}

Durante las pruebas de campo se observó que la ruta ``saltaba'' a otros lugares. El análisis de
\textbf{24\,351 posiciones reales} mostró que todos esos saltos coincidían con lecturas que el propio
teléfono reportaba con un margen de error de 74 a 800\,m: ubicaciones aproximadas por Wi-Fi o antenas
celulares, no por satélite. La Figura~\ref{fig:precision} muestra la distribución.

\begin{figure}[!htbp]
\centering
\begin{tikzpicture}
\begin{axis}[
  ybar, bar shift=0pt, width=\linewidth, height=6.2cm, bar width=0.9cm,
  ylabel={\% de lecturas}, xlabel={Margen de error reportado por el GPS (m)},
  symbolic x coords={0--5,5--10,10--15,15--20,20--25,25--40,40--100,>100},
  xtick={0--5,5--10,10--15,15--20,20--25,25--40,40--100,>100}, ymin=0, ymax=62,
  enlarge x limits=0.07,
  nodes near coords={\pgfmathprintnumber[fixed,precision=1]{\pgfplotspointmeta}\%},
  nodes near coords style={font=\scriptsize, text=inkgray},
  every axis plot/.append style={draw=none},
  axis lines*=left, ymajorgrids, grid style={gray!20},
  tick label style={font=\footnotesize}, label style={font=\small},
]
\addplot[fill=ggreen!70] coordinates {(0--5,15.1) (5--10,54.1) (10--15,14.3) (15--20,5.7) (20--25,3.1)};
\addplot[fill=gyellow!80] coordinates {(25--40,4.0)};
\addplot[fill=gred!75] coordinates {(40--100,2.8) (>100,0.9)};
\node[font=\scriptsize, text=gred, align=center] at (axis cs:40--100,24) {lecturas que\\``teletransportan''\\\faArrowDown};
\end{axis}
\end{tikzpicture}
\caption{Distribución de la precisión en 24\,351 lecturas reales. En verde: aceptadas siempre; amarillo:
aceptadas solo tras 20\,s sin una lectura mejor; rojo: descartadas.}
\label{fig:precision}
\end{figure}

Para resolverlo se diseñó un filtro de cuatro etapas que procesa cada lectura antes de guardarla, enviarla o
dibujarla (Figura~\ref{fig:filtro}).

\begin{figure}[!htbp]
\centering
\begin{tikzpicture}[node distance=0.55cm and 0.9cm, font=\footnotesize]
  \node[box=gblue, minimum width=3cm] (in) {\faSatelliteDish\ Lectura GPS\\{\scriptsize lat, lng, $\sigma$, $v$, rumbo}};
  \node[decision, below=of in] (d1) {¿$\sigma \le 25$\,m?\\{\scriptsize (40\,m tras 20\,s)}};
  \node[decision, below=of d1] (d2) {¿velocidad\\implícita $\le 50$\,m/s?};
  \node[decision, below=of d2] (d3) {¿pausa\\$> 10$\,s?};
  \node[box=ggreen, below=of d3, minimum width=3.4cm] (kal) {\textbf{Kalman}\\{\scriptsize predicción con $v$ y rumbo}\\{\scriptsize + corrección ponderada}};
  \node[decision, below=of kal] (d4) {¿quieto\\3 lecturas?};
  \node[box=purple, below left=0.6cm and -0.2cm of d4, minimum width=3.2cm] (avg) {\textbf{Posición fija}\\{\scriptsize promedio ponderado $1/\sigma^2$}};
  \node[box=ggreen, below right=0.6cm and -0.2cm of d4, minimum width=3.2cm] (mov) {\textbf{En movimiento}\\{\scriptsize salida del Kalman}};
  \node[box=inkgray, below=2.6cm of d4, minimum width=4.2cm] (out) {\faMapMarker*\ Posición filtrada\\{\scriptsize SQLite, servidor y mapa}};

  \node[box=gred, right=1.6cm of d1, minimum width=2.4cm] (x1) {\faTrash\ Descartar};
  \node[box=gred, right=1.6cm of d2, minimum width=2.4cm] (x2) {\faTrash\ Descartar\\{\scriptsize 5 seguidas \faArrowRight\ reiniciar}};
  \node[box=gyellow!90!black, left=1.4cm of d3, minimum width=2.4cm] (rst) {\faRedo\ Reiniciar\\{\scriptsize en la lectura nueva}};

  \draw[arr] (in) -- (d1);
  \draw[arr] (d1) -- node[lbl, right]{sí} (d2);
  \draw[arr] (d1) -- node[lbl, above]{no} (x1);
  \draw[arr] (d2) -- node[lbl, right]{sí} (d3);
  \draw[arr] (d2) -- node[lbl, above]{no} (x2);
  \draw[arr] (d3) -- node[lbl, right]{no} (kal);
  \draw[arr] (d3) -- node[lbl, above]{sí} (rst);
  \draw[arr] (rst) |- (out.west);
  \draw[arr] (kal) -- (d4);
  \draw[arr] (d4) -| node[lbl, pos=0.25, above]{sí} (avg);
  \draw[arr] (d4) -| node[lbl, pos=0.25, above]{no} (mov);
  \draw[arr] (avg) |- (out);
  \draw[arr] (mov) |- (out);
\end{tikzpicture}
\caption{Etapas del filtro de precisión (\texttt{location\_filter.dart}).}
\label{fig:filtro}
\end{figure}

\paragraph{Filtro de Kalman con predicción.}
Para cada coordenada se mantiene una estimación $\hat{x}$ y su varianza $P$ (en m$^2$). Con un intervalo
$\Delta t$, una velocidad $v$ y un rumbo $\theta$ reportados por el GPS (efecto Doppler), y una medición $z$
de precisión $\sigma$:
\begin{align}
  \hat{x}^- &= \hat{x} + v\,\Delta t\,(\cos\theta,\ \sin\theta) && \text{predicción por velocidad y rumbo}\\
  P^- &= P + \Delta t\, q^2, \qquad q = 2 + 0.5\,v && \text{incertidumbre del movimiento}\\
  K &= \frac{P^-}{P^- + \sigma^2} && \text{ganancia: cuánto confiar en la medición}\\
  \hat{x} &= \hat{x}^- + K\,(z - \hat{x}^-), \qquad P = (1-K)\,P^- && \text{corrección}
\end{align}
Predecir con la velocidad evita que el punto ``se quede atrás'' en un vehículo, y hacer que $q$ crezca con
la velocidad permite seguir curvas y frenadas.

\paragraph{Posición fija estando quieto.}
Cuando la velocidad es menor a 0.5\,m/s durante tres lecturas, la posición se fija en un promedio ponderado
con olvido gradual ($\lambda = 0.98$, memoria de unos 50\,s), que da más peso a las lecturas precisas:
\begin{equation}
  \bar{x} = \frac{\sum_i \lambda^{n-i}\, w_i\, z_i}{\sum_i \lambda^{n-i}\, w_i}, \qquad w_i = \frac{1}{\sigma_i^2}
\end{equation}
El modo se abandona cuando la velocidad supera 1\,m/s o dos lecturas seguidas caen lejos del punto fijo.

\begin{lstlisting}[caption={Núcleo del filtro de Kalman (Dart).}]
// Predicción: avanzar según la velocidad y rumbo que reporta el GPS (Doppler).
if (speed != null && speed > 0.5 && heading != null && heading >= 0) {
  final step = speed * dt;
  final h = heading * pi / 180;
  _lat = _lat! + (step * cos(h)) / 111320.0;
  _lng = _lng! + (step * sin(h)) / (111320.0 * cos(_lat! * pi / 180));
}
final q = processNoise + 0.5 * max(speed ?? 0, 0);
_variance += dt * q * q;

// Corrección con la medición, ponderada por su precisión.
final k = _variance / (_variance + acc * acc);
_lat = _lat! + k * (lat - _lat!);
_lng = _lng! + k * (lng - _lng!);
_variance = (1 - k) * _variance;
\end{lstlisting}

\subsection{Telemetría en el dispositivo}

La velocidad se toma del GPS (efecto Doppler, más estable que distancia/tiempo) y se suaviza con la mediana
de las tres últimas lecturas. Con ella se clasifica la actividad y se estiman las calorías:
\begin{equation}
  \text{kcal} = \text{MET}(\text{actividad}, v) \times \text{peso (kg)} \times \text{tiempo (h)}
\end{equation}

\begin{figure}[!htbp]
\centering
\begin{tikzpicture}[x=0.72cm]
  % Escala no lineal para que cada franja tenga espacio para su etiqueta
  \fill[gray!35]      (0,0) rectangle (3.5,0.8);
  \fill[ggreen!60]    (3.5,0) rectangle (8.5,0.8);
  \fill[gyellow!75]   (8.5,0) rectangle (14,0.8);
  \fill[gred!55]      (14,0) rectangle (19.5,0.8);
  \node[font=\footnotesize\bfseries] at (1.75,0.4) {\faStopCircle\ Detenido};
  \node[font=\footnotesize\bfseries] at (6,0.4) {\faWalking\ Caminata};
  \node[font=\footnotesize\bfseries] at (11.25,0.4) {\faRunning\ Alta velocidad};
  \node[font=\footnotesize\bfseries] at (16.75,0.4) {\faCar\ Vehículo};
  \foreach \v/\k in {0/0, 3.5/0.5, 8.5/2.0, 14/6.0}{
    \draw[thick] (\v,0) -- (\v,-0.15) node[below, font=\scriptsize]{\k};
  }
  \node[below, font=\scriptsize] at (19.5,-0.15) {m/s};
  \node[font=\scriptsize, text=gray, align=center] at (1.75,-0.95) {0 kcal};
  \node[font=\scriptsize, text=gray, align=center] at (6,-0.95) {MET 2.0--5.0\\según ritmo};
  \node[font=\scriptsize, text=gray, align=center] at (11.25,-0.95) {MET $\approx$ km/h\\(mín. 6.0)};
  \node[font=\scriptsize, text=gray, align=center] at (16.75,-0.95) {no suma calorías\\(más de 21.6 km/h)};
\end{tikzpicture}
\caption{Clasificación de la actividad por velocidad y MET asociado (Compendium of Physical Activities).}
\end{figure}

Para no inflar la distancia, los desplazamientos solo se suman si la lectura es precisa ($\sigma \le 35$\,m)
y hay movimiento real (velocidad $> 0.5$\,m/s o desplazamiento mayor que el margen de error).

\subsection{Interfaz móvil}

La interfaz sigue la estética de Google Maps: el mapa ocupa toda la pantalla y los controles flotan encima
en tarjetas blancas con sombra suave.

\begin{figure}[H]
\centering
\capturaO{app_mapa.png}{0.32\linewidth}{%
\begin{tikzpicture}[scale=0.9, every node/.style={font=\tiny}]
  % Teléfono
  \draw[rounded corners=14pt, fill=inkgray, draw=inkgray] (-0.25,-0.3) rectangle (5.25,10.3);
  \begin{scope}
    \clip[rounded corners=10pt] (0,0) rectangle (5,10);
    \fill[mapbg] (0,0) rectangle (5,10);
    \foreach \y in {1.3,3.1,5.4,7.6,9.1} \draw[white, line width=3pt] (0,\y) -- (5,\y+0.4);
    \foreach \x in {0.9,2.4,3.9} \draw[white, line width=3pt] (\x,0) -- (\x+0.3,10);
    \fill[gblue!12] (3.2,7.9) rectangle (4.6,8.8);
    \fill[ggreen!25] (0.2,1.6) rectangle (1.5,2.8);
    % Ruta
    \draw[white, line width=4pt, rounded corners=3pt] (1.0,2.0) -- (1.1,3.3) -- (2.5,3.6) -- (2.6,5.6);
    \draw[gblue, line width=2.2pt, rounded corners=3pt] (1.0,2.0) -- (1.1,3.3) -- (2.5,3.6) -- (2.6,5.6);
    % Precisión + punto
    \fill[gblue, opacity=0.15] (2.6,5.6) circle (0.55);
    \draw[gblue, opacity=0.4] (2.6,5.6) circle (0.55);
    \fill[white] (2.6,5.6) circle (0.17);
    \fill[gblue] (2.6,5.6) circle (0.12);
    % Barra superior
    \fill[white, rounded corners=8pt, drop shadow={opacity=0.2}] (0.3,9.0) rectangle (4.7,9.7);
    \node[anchor=west, text=inkgray] at (0.5,9.35) {\textbf{GeoTrack}};
    \node[anchor=east, text=gray] at (4.5,9.35) {\faHistory\ \ \faCog};
    % Botones laterales
    \fill[white, drop shadow={opacity=0.2}] (4.4,6.4) circle (0.28); \node[text=gray] at (4.4,6.4) {\faCrosshairs};
    \fill[gred, drop shadow={opacity=0.2}] (4.4,5.6) circle (0.28); \node[text=white] at (4.4,5.6) {\textbf{SOS}};
    % Tarjeta inferior
    \fill[white, rounded corners=8pt, drop shadow={opacity=0.2}] (0.3,0.3) rectangle (4.7,2.1);
    \node[anchor=west, text=inkgray] at (0.5,1.85) {\faWalking\ \textbf{Caminata}};
    \node[anchor=west, text=ggreen] at (0.5,1.5) {\faCircle\ \textcolor{gray}{Sincronizado · GPS $\pm$5 m}};
    \node[text=inkgray, align=center] at (0.95,0.85) {\textbf{4.8}\\km/h};
    \node[text=inkgray, align=center] at (2.0,0.85) {\textbf{1.27}\\km};
    \node[text=inkgray, align=center] at (3.0,0.85) {\textbf{86}\\kcal};
    \fill[gred, drop shadow={opacity=0.3}] (4.1,0.9) circle (0.38); \node[text=white, font=\small] at (4.1,0.9) {\faStop};
  \end{scope}
\end{tikzpicture}}
\hspace{0.8cm}
\begin{minipage}[b]{0.55\linewidth}\small
\begin{itemize}
  \item \textbf{Ruta en vivo} dibujada con borde blanco, como en Google Maps.
  \item \textbf{Punto azul con círculo de precisión}: el radio es el margen de error del GPS en ese momento.
  \item \textbf{Línea de estado}: rastreo, sincronización, puntos en cola y precisión (``GPS $\pm$5\,m'').
  \item \textbf{Tarjeta de telemetría}: velocidad, distancia, calorías y actividad.
  \item \textbf{Botón SOS}: se mantiene presionado 2\,s para evitar activaciones accidentales.
  \item \textbf{Botón flotante} Iniciar/Detener.
\end{itemize}
\vspace{1.5cm}
\end{minipage}
\caption{Pantalla principal de la aplicación móvil.}
\end{figure}

% ════════════════════════════════════════════════════════════════════════════
\section{Backend: API y tiempo real}

\subsection{Modelo de datos}

\begin{figure}[H]
\centering
\begin{tikzpicture}[
  ent/.style={draw=#1, fill=white, rounded corners=3pt, thick, inner sep=1pt, font=\scriptsize,
    drop shadow={opacity=0.12}},
  rel/.style={thick, color=inkgray, -{Stealth[length=2mm]}},
  node distance=1.0cm and 1.3cm]
  \node[ent=gblue] (dev) {\entidad{gblue}{devices}{\faKey\ id\\name, color, platform\\last\_seen\_at\\last\_lat, last\_lng\\last\_speed, last\_battery\\tracking}};
  \node[ent=ggreen, right=of dev] (ses) {\entidad{ggreen}{sessions}{\faKey\ id\\device\_id\\name\\started\_at, ended\_at\\distance\_m, calories}};
  \node[ent=awsorange, right=of ses] (loc) {\entidad{awsorange}{locations}{\faKey\ id\\device\_id, session\_id\\ts {\tiny (UNIQUE con device\_id)}\\lat, lng, accuracy\\speed, heading, altitude\\battery, activity}};
  \node[ent=gred, below=of dev] (alr) {\entidad{gred}{alerts}{\faKey\ id\\device\_id, type\\message, lat, lng\\created\_at\\acknowledged}};
  \node[ent=purple, below=of ses] (gst) {\entidad{purple}{geofence\_state}{\faKey\ device\_id\\\faKey\ geofence\_id\\inside}};
  \node[ent=purple, below=of loc] (geo) {\entidad{purple}{geofences}{\faKey\ id\\name\\lat, lng, radius\_m}};
  \draw[rel] (ses) -- node[lbl, above]{N:1} (dev);
  \draw[rel] (loc) -- node[lbl, above]{N:1} (ses);
  \draw[rel] (loc.north) -- ++(0,0.5) -| node[lbl, pos=0.25, above]{N:1} (dev.north);
  \draw[rel] (alr) -- node[lbl, right]{N:1} (dev);
  \draw[rel] (gst) -- node[lbl, above]{N:1} (geo);
  \draw[rel] (gst.north west) -- node[lbl, pos=0.4, right=2pt]{N:1} (dev.south east);
\end{tikzpicture}
\caption{Modelo entidad-relación en PostgreSQL.}
\end{figure}

\subsection{API REST}

\begin{table}[H]
\centering
\small
\begin{tabularx}{\linewidth}{l l X}
\toprule
\textbf{Método} & \textbf{Ruta} & \textbf{Descripción} \\
\midrule
\multicolumn{3}{l}{\textcolor{gblue}{\faMobile*\ \textbf{Dispositivos}} (cabecera \texttt{x-api-key})} \\
POST & \texttt{/api/devices/register} & Registra el teléfono y le asigna un color único. \\
POST & \texttt{/api/sessions/start} & Inicia una sesión de rastreo. \\
POST & \texttt{/api/sessions/:id/stop} & Cierra la sesión con distancia y calorías finales. \\
PATCH & \texttt{/api/sessions/:id} & Renombra una sesión. \\
POST & \texttt{/api/locations} & Recibe un lote de posiciones (idempotente). \\
POST & \texttt{/api/alerts/sos} & Emite una alerta SOS con la última posición. \\
\midrule
\multicolumn{3}{l}{\textcolor{ggreen}{\faDesktop\ \textbf{Dashboard}}} \\
GET & \texttt{/api/devices}, \texttt{/api/live} & Estado de los dispositivos y rutas de las sesiones activas. \\
GET & \texttt{/api/history} & Posiciones de un dispositivo en un rango de fechas. \\
GET & \texttt{/api/history/stats} & Distancia, duración, velocidad media y máxima. \\
GET & \texttt{/api/history/export} & Descarga del recorrido en GPX o CSV. \\
GET/POST/DELETE & \texttt{/api/geofences} & Administración de geocercas. \\
GET/POST & \texttt{/api/alerts}, \texttt{/:id/ack} & Listado y confirmación de alertas. \\
\bottomrule
\end{tabularx}
\caption{Principales rutas de la API.}
\end{table}

\subsection{Procesamiento de cada lote}

\begin{figure}[!htbp]
\centering
\begin{tikzpicture}[node distance=0.4cm, font=\footnotesize]
  \node[box=gblue, minimum width=2.4cm] (a) {\faInbox\\Lote recibido};
  \node[box=awsorange, right=of a, minimum width=2.4cm] (b) {\faDatabase\\INSERT\\idempotente};
  \node[box=awsorange, right=of b, minimum width=2.4cm] (c) {\faSync\\Actualizar\\último estado};
  \node[box=gred, right=of c, minimum width=2.4cm] (d) {\faBatteryQuarter\\¿Batería\\$\le 15$\,\%?};
  \node[box=purple, right=of d, minimum width=2.4cm] (e) {\faDrawPolygon\\Evaluar\\geocercas};
  \node[box=ggreen, below=0.9cm of c, minimum width=5cm] (f) {\faBroadcastTower\ \texttt{emit('locations')} + \texttt{emit('device')}};
  \draw[arr] (a) -- (b); \draw[arr] (b) -- (c); \draw[arr] (c) -- (d); \draw[arr] (d) -- (e);
  \draw[arr] (e) |- (f);
\end{tikzpicture}
\caption{Flujo de \texttt{POST /api/locations} en el servidor.}
\end{figure}

Además, un \textbf{vigilante} se ejecuta cada 30\,s y genera una alerta \emph{``sin señal''} si un
dispositivo que está rastreando lleva más de 2 minutos sin enviar datos.

\subsection{Geocercas}
Una geocerca es un círculo (centro y radio). Para cada posición se calcula la distancia al centro con la
fórmula de Haversine, y se compara con el último estado conocido (dentro/fuera) guardado en
\texttt{geofence\_state}; solo un \emph{cambio} de estado genera alerta:
\begin{equation}
  d = 2R \arcsin\sqrt{\sin^2\!\frac{\Delta\varphi}{2} + \cos\varphi_1\cos\varphi_2\sin^2\!\frac{\Delta\lambda}{2}},
  \qquad R = 6\,371\,000\ \text{m}
\end{equation}

\begin{figure}[H]
\centering
\begin{tikzpicture}[node distance=3.4cm, font=\footnotesize]
  \node[circle, draw=gray, fill=softgray, thick, minimum size=1.6cm, align=center] (u) {Sin\\estado};
  \node[circle, draw=ggreen, fill=ggreen!12, thick, minimum size=1.6cm, right=of u] (in) {Dentro};
  \node[circle, draw=gred, fill=gred!10, thick, minimum size=1.6cm, right=of in] (out) {Fuera};
  \draw[arr] (u) -- node[lbl, above]{1.\textsuperscript{a} lectura}node[lbl, below]{sin alerta} (in);
  \draw[arr] (u.north) to[bend left=35] node[lbl, above]{1.\textsuperscript{a} lectura, sin alerta} (out.north);
  \draw[arr, gred] (in) to[bend left=20] node[lbl, above]{$d > r$ \faBell\ \textbf{salió}} (out);
  \draw[arr, ggreen] (out) to[bend left=20] node[lbl, below]{$d \le r$ \faBell\ \textbf{entró}} (in);
\end{tikzpicture}
\caption{Máquina de estados de una geocerca por dispositivo.}
\end{figure}

% ════════════════════════════════════════════════════════════════════════════
\section{Centro de monitoreo web}

\begin{figure}[H]
\centering
\capturaO{dashboard.png}{\linewidth}{%
\begin{tikzpicture}[every node/.style={font=\tiny}]
  \fill[rounded corners=4pt, softgray, draw=gray!50] (0,0) rectangle (15,8.4);
  \fill[gray!25, rounded corners=4pt] (0,7.9) rectangle (15,8.4);
  \foreach \x/\c in {0.3/gred,0.6/gyellow,0.9/ggreen} \fill[\c] (\x,8.15) circle (0.08);
  \begin{scope}
    \clip (0,0) rectangle (15,7.9);
    \fill[mapbg] (0,0) rectangle (15,7.9);
    \foreach \y in {1.2,2.9,4.6,6.4} \draw[white, line width=4pt] (0,\y) -- (15,\y+0.6);
    \foreach \x in {5.5,8.2,11,13.4} \draw[white, line width=4pt] (\x,0) -- (\x-0.5,7.9);
    \fill[gblue!12] (9,5.1) rectangle (10.6,6.0);
    \fill[ggreen!25] (11.6,1.6) rectangle (13,2.6);
    % Rutas
    \foreach \c/\p in {gblue/{(6,1.5) (6.4,3.2) (8.0,3.6) (8.4,5.8) (10.2,6.3)}, gred/{(13.6,0.8) (13.2,2.2) (11.2,2.0) (11.0,4.2) (12.4,4.6)}, purple/{(4.8,6.8) (5.3,5.2) (7.0,4.9) (7.6,4.0)}}{
      \draw[white, line width=4pt, rounded corners=3pt] plot coordinates {\p};
      \draw[\c, line width=2pt, rounded corners=3pt] plot coordinates {\p};
    }
    \foreach \c/\x/\y/\r in {gblue/10.2/6.3/0.35, gred/12.4/4.6/0.25, purple/7.6/4.0/0.3}{
      \fill[\c, opacity=0.15] (\x,\y) circle (\r);
      \fill[white] (\x,\y) circle (0.15); \fill[\c] (\x,\y) circle (0.1);
    }
    \node[fill=white, rounded corners=2pt, inner sep=2pt] at (10.2,6.75) {\textcolor{gblue}{\faCircle} Mi teléfono · 4.8 km/h};
    \node[fill=white, rounded corners=2pt, inner sep=2pt] at (12.4,5.0) {\textcolor{gred}{\faCircle} Repartidor 2 · 32 km/h};
    \node[fill=white, rounded corners=2pt, inner sep=2pt] at (7.6,4.4) {\textcolor{purple}{\faCircle} Supervisor · detenido};
  \end{scope}
  % Panel lateral
  \fill[white, rounded corners=6pt, drop shadow={opacity=0.2}] (0.3,1.4) rectangle (4.2,7.6);
  \node[anchor=west, text=inkgray, font=\scriptsize\bfseries] at (0.5,7.25) {\faMapMarker*\ GeoTrack · Monitoreo};
  \foreach \i/\t in {0/\faBroadcastTower, 1/\faHistory, 2/\faDrawPolygon, 3/\faBell}{
    \node[anchor=center, text=gray, font=\scriptsize] at (0.95+\i*0.9,6.75) {\t};
  }
  \node[anchor=west, text=gblue, font=\tiny\bfseries] at (0.5,6.42) {En vivo};
  \fill[gblue] (0.6,6.56) rectangle (1.3,6.6);
  \foreach \y/\c/\n/\s in {5.8/gblue/Mi teléfono/Caminata · 86\,\%, 5.0/gred/Repartidor 2/Vehículo · 54\,\%, 4.2/purple/Supervisor/Detenido · 71\,\%}{
    \fill[softgray, rounded corners=3pt] (0.5,\y-0.32) rectangle (4.0,\y+0.32);
    \fill[\c] (0.75,\y) circle (0.12);
    \node[anchor=west, text=inkgray] at (0.95,\y+0.1) {\textbf{\n}};
    \node[anchor=west, text=gray] at (0.95,\y-0.15) {\s};
    \node[anchor=east, text=ggreen, font=\tiny] at (3.9,\y) {\faCircle};
  }
  % Reproductor
  \fill[white, rounded corners=10pt, drop shadow={opacity=0.2}] (5.2,0.3) rectangle (14.4,1.15);
  \fill[gblue] (5.65,0.72) circle (0.25); \node[text=white] at (5.67,0.72) {\faPlay};
  \draw[gray!40, line width=2.5pt, line cap=round] (6.2,0.72) -- (11.4,0.72);
  \draw[gblue, line width=2.5pt, line cap=round] (6.2,0.72) -- (8.6,0.72);
  \fill[gblue] (8.6,0.72) circle (0.1);
  \foreach \i/\v in {0/x1,1/x10,3/x30,4/x60}{
    \node[draw=gray!50, rounded corners=5pt, inner sep=2pt] at (11.85+\i*0.55,0.72) {\v};
  }
  \node[draw=gblue, fill=gblue, text=white, rounded corners=5pt, inner sep=2pt] at (11.85+2*0.55,0.72) {x20};
\end{tikzpicture}}
\caption{Centro de monitoreo: panel en vivo con varios dispositivos y reproductor de rutas.}
\end{figure}

\subsection{Funcionalidades}
\begin{itemize}
  \item \textbf{Panel en vivo:} todos los dispositivos que rastrean, cada uno con su color, etiqueta con nombre
        y velocidad, y círculo de precisión alrededor de la última posición.
  \item \textbf{Trazado en tiempo real:} cada evento \texttt{locations} de Socket.IO extiende la polilínea
        del dispositivo sin recargar la página.
  \item \textbf{Consultas históricas:} filtros por dispositivo, fecha y hora, con atajos (1\,h, Hoy, 24\,h,
        7 días) y estadísticas del período.
  \item \textbf{Reproductor (playback):} Play/Pausa, barra de progreso arrastrable y velocidades x1, x10,
        x20, x30 y x60; los huecos sin señal se saltan para no esperar en vano.
  \item \textbf{Exportación} del recorrido en GPX (compatible con Google Earth y Strava) o CSV.
\end{itemize}

% ════════════════════════════════════════════════════════════════════════════
\section{Innovaciones (Fase 5)}

\begin{table}[H]
\centering
\small
\begin{tabularx}{\linewidth}{>{\centering\arraybackslash}p{0.8cm} l X}
\toprule
 & \textbf{Funcionalidad} & \textbf{Valor que aporta} \\
\midrule
\textcolor{purple}{\Large\faDrawPolygon} & Geocercas con alertas & Avisa cuando un dispositivo entra o sale de una zona (almacén, colegio, obra). Se dibujan con un clic y un radio. \\
\textcolor{gred}{\Large\faExclamationCircle} & Botón SOS & Alerta roja inmediata en el panel con la posición exacta; se mantiene 2\,s para evitar falsas alarmas. \\
\textcolor{gyellow!90!black}{\Large\faBatteryQuarter} & Alertas automáticas & Batería baja ($\le 15$\,\%) y dispositivo \emph{sin señal} ($> 2$\,min sin datos). \\
\textcolor{ggreen}{\Large\faFileExport} & Exportación GPX/CSV y estadísticas & Distancia, duración, velocidad media y máxima; integración con otras herramientas. \\
\textcolor{gblue}{\Large\faLock} & Autenticación & API key para dispositivos y token para el panel. \\
\textcolor{inkgray}{\Large\faCrosshairs} & Filtro de precisión GNSS & Elimina saltos, suaviza la ruta, fija la posición estando quieto y muestra el margen de error. \\
\bottomrule
\end{tabularx}
\caption{Funcionalidades adicionales implementadas.}
\end{table}

% ════════════════════════════════════════════════════════════════════════════
\section{Despliegue y operación}

\subsection{Scripts de infraestructura}

\begin{table}[H]
\centering
\small
\begin{tabularx}{\linewidth}{l X}
\toprule
\textbf{Script} & \textbf{Qué hace} \\
\midrule
\texttt{deploy-aws.ps1} & Crea llave SSH, Security Groups, RDS, EC2 con Nginx y Node.js, IP elástica, y despliega la aplicación (unos 15 min). \\
\texttt{update-app.ps1} & Compila el dashboard, empaqueta el backend, lo sube por SSH, reinicia el servicio y verifica \texttt{/api/health}. \\
\texttt{build-apk-cloud.ps1} & Alternativa: compila el APK en una EC2 Linux temporal que se autodestruye. \\
\texttt{destroy-aws.ps1} & Elimina todos los recursos para dejar de pagar. \\
\bottomrule
\end{tabularx}
\caption{Automatización del ciclo de vida de la infraestructura.}
\end{table}

\begin{figure}[!htbp]
\centering
\begin{tikzpicture}[node distance=0.35cm, font=\footnotesize]
  \node[box=inkgray, minimum width=2.2cm] (code) {\faCode\\Cambio de\\código};
  \node[box=gblue, right=of code, minimum width=2.2cm] (an) {\faCheckDouble\\\texttt{dart analyze}};
  \node[box=gblue, right=of an, minimum width=2.4cm] (apk) {\faCogs\\\texttt{flutter build}\\\texttt{apk --release}};
  \node[box=gblue, right=of apk, minimum width=2.2cm] (adb) {\faMobile*\\\texttt{adb install -r}};
  \node[box=ggreen, below=0.9cm of an, minimum width=2.2cm] (vite) {\faCubes\\\texttt{vite build}};
  \node[box=awsorange, right=of vite, minimum width=2.4cm] (scp) {\faUpload\\\texttt{scp} + \texttt{npm ci}};
  \node[box=awsorange, right=of scp, minimum width=2.2cm] (sysd) {\faRedo\\\texttt{systemctl}\\\texttt{restart}};
  \node[box=ggreen, right=of sysd, minimum width=2.2cm] (hc) {\faHeartbeat\\\texttt{/api/health}};
  \draw[arr] (code) -- (an); \draw[arr] (an) -- (apk); \draw[arr] (apk) -- (adb);
  \draw[arr] (code) |- (vite); \draw[arr] (vite) -- (scp); \draw[arr] (scp) -- (sysd); \draw[arr] (sysd) -- (hc);
  \node[font=\scriptsize\bfseries, text=gblue, above=0.05cm of apk] {App móvil};
  \node[font=\scriptsize\bfseries, text=awsorange!80!black, below=0.05cm of scp] {Servidor y dashboard (\texttt{update-app.ps1})};
\end{tikzpicture}
\caption{Flujo de entrega de una nueva versión.}
\end{figure}

\begin{nota}[Firma del APK]
El APK de producción se firma siempre con la misma llave (\texttt{geotrack-release.jks}). Así cada
actualización se instala encima de la anterior sin borrar el historial guardado en el teléfono.
\end{nota}

\subsection{Problemas encontrados y soluciones}

\begin{table}[H]
\centering
\small
\begin{tabularx}{\linewidth}{>{\raggedright\arraybackslash}p{3.6cm} X X}
\toprule
\textbf{Problema} & \textbf{Causa} & \textbf{Solución} \\
\midrule
La ruta ``salta'' a otros lugares & Lecturas por Wi-Fi/antenas con error de 74--800\,m tratadas como exactas. & Filtro de precisión de cuatro etapas (Sección 5.4). \\
El punto se mueve estando quieto & Ruido normal del GNSS de 5--20\,m. & Kalman + promedio ponderado en modo quieto. \\
Aviso falso ``Sin conexión'' & La pantalla solo se refrescaba al llegar una posición nueva. & Refresco tras cada envío y estado inicial ``desconocido''. \\
No compila en Windows & La tilde de la carpeta (``posición'') rompe el compilador AOT de Dart. & Unidad virtual sin tildes (\texttt{subst G:}) y caché incremental de Kotlin desactivada. \\
Scripts colgados en SSH & Cambió la IP pública del administrador. & Nueva regla \texttt{/32} en el Security Group. \\
Android pausa el GPS & Optimización de batería del fabricante. & Servicio en primer plano + watchdog cada 30\,s. \\
\bottomrule
\end{tabularx}
\caption{Incidencias técnicas y su resolución.}
\end{table}

% ════════════════════════════════════════════════════════════════════════════
\section{Resultados y pruebas de campo}

El sistema se probó durante varios días en la ciudad del Cusco con un Samsung Galaxy A26 (receptor GNSS
Samsung S.LSI), caminando, en vehículo y dentro de edificios.

\begin{table}[H]
\centering
\small
\begin{tabular}{l r}
\toprule
\textbf{Indicador} & \textbf{Valor} \\
\midrule
Dispositivos registrados & 3 \\
Sesiones de rastreo & 21 \\
Posiciones almacenadas en el servidor & 50\,785 \\
Distancia acumulada en sesiones & 97.8 km \\
Velocidad máxima registrada (vehículo) & 71.4 km/h \\
Lecturas con precisión $\le 10$\,m & 69.2\,\% \\
Intervalo mediano entre puntos (configurado a 1\,s) & 1.0 s \\
\bottomrule
\end{tabular}
\caption{Resumen de las pruebas de campo.}
\end{table}

\subsection{Evaluación del filtro de precisión}

Las 24\,351 lecturas crudas del 8 y 9 de octubre se volvieron a procesar con el filtro para medir su
efecto de forma objetiva:

\begin{figure}[!htbp]
\centering
\begin{subfigure}[b]{0.48\linewidth}
\centering
\begin{tikzpicture}
\begin{axis}[
  ybar, bar shift=0pt, width=\linewidth, height=5.4cm, bar width=0.75cm, ymin=0, ymax=23,
  symbolic x coords={Crudo, Kalman, {Kalman + quieto}}, xtick={Crudo, Kalman, {Kalman + quieto}},
  ylabel={km}, nodes near coords, nodes near coords style={font=\scriptsize},
  axis lines*=left, ymajorgrids, grid style={gray!20}, tick label style={font=\scriptsize},
  every axis plot/.append style={draw=none}, enlarge x limits=0.25,
]
\addplot[fill=gred!70] coordinates {(Crudo,18.29)};
\addplot[fill=gyellow!80] coordinates {(Kalman,6.37)};
\addplot[fill=ggreen!70] coordinates {({Kalman + quieto},3.59)};
\end{axis}
\end{tikzpicture}
\caption{Movimiento falso acumulado estando quieto.}
\end{subfigure}\hfill
\begin{subfigure}[b]{0.48\linewidth}
\centering
\begin{tikzpicture}
\begin{axis}[
  ybar, bar shift=0pt, width=\linewidth, height=5.4cm, bar width=0.75cm, ymin=0, ymax=11,
  symbolic x coords={Crudo, Filtrado}, xtick={Crudo, Filtrado},
  ylabel={saltos detectados}, nodes near coords, nodes near coords style={font=\scriptsize},
  axis lines*=left, ymajorgrids, grid style={gray!20}, tick label style={font=\scriptsize},
  every axis plot/.append style={draw=none}, enlarge x limits=0.45,
]
\addplot[fill=gred!70] coordinates {(Crudo,9)};
\addplot[fill=ggreen!70] coordinates {(Filtrado,2)};
\end{axis}
\end{tikzpicture}
\caption{Saltos a otro lugar y regreso.}
\end{subfigure}
\caption{Efecto del filtro sobre datos reales.}
\end{figure}

\begin{table}[H]
\centering
\small
\begin{tabular}{l c c}
\toprule
\textbf{Métrica} & \textbf{Sin filtro} & \textbf{Con filtro} \\
\midrule
Saltos a otro lugar & 9 & 2\textsuperscript{*} \\
Movimiento falso estando quieto & 18.29 km & 3.59 km ($-80$\,\%) \\
Desfase respecto al GPS en vehículo (percentil 95) & — & 5.9 m \\
Lecturas descartadas por imprecisas & 0 \% & 6.8 \% \\
\bottomrule
\end{tabular}
\caption{Métricas del filtro. \textsuperscript{*}Los dos restantes corresponden a vueltas en U reales
después de que el GPS se pausara.}
\end{table}

\begin{alerta}[Límite físico]
Un receptor GNSS de teléfono de una sola frecuencia tiene un límite de unos 3--5\,m al aire libre y peor
entre edificios o bajo techo, por la reflexión de la señal (\emph{multipath}). Ningún software lo elimina por
completo; para precisión centimétrica se requieren equipos con corrección RTK. Por eso la app muestra el
\textbf{círculo de precisión}: comunica honestamente cuánto confiar en cada lectura.
\end{alerta}

% ════════════════════════════════════════════════════════════════════════════
\section{Conclusiones}

\begin{enumerate}
  \item Se construyó y desplegó un sistema de rastreo completo y funcional de extremo a extremo: app
        Android, backend en AWS y centro de monitoreo web, que cumple los diez requerimientos funcionales y
        supera el mínimo de innovaciones solicitado.
  \item El enfoque \emph{offline-first} con idempotencia en el servidor permite que los cortes de red no
        pierdan ni dupliquen datos: todo se guarda primero en el teléfono y se reenvía al volver la señal, algo esencial en un sistema embebido móvil con conectividad intermitente.
  \item El servicio en primer plano con watchdog, independiente de la interfaz, permitió un rastreo continuo
        de un punto por segundo (intervalo mediano medido en campo: 1.0\,s).
  \item El análisis de datos reales permitió diagnosticar el origen de los errores de posición y diseñar un
        filtro (validación, Kalman y promedio estando quieto) que eliminó los saltos y redujo un 80\,\% la
        deriva estando quieto.
  \item La automatización con scripts hizo la infraestructura reproducible y redujo cada actualización del
        servidor a un solo comando.
\end{enumerate}

\subsection{Trabajo futuro}
\begin{itemize}
  \item \textbf{HTTPS} con un dominio propio y certificado de Let's Encrypt (\texttt{certbot}).
  \item \textbf{Ajuste a calles} (\emph{map matching}) para pegar la ruta del vehículo a la vía.
  \item Kalman de estado completo (posición y velocidad) y fusión con acelerómetro y giroscopio.
  \item Notificaciones \emph{push} de las alertas al teléfono del supervisor.
  \item Cuentas de usuario con roles y organizaciones (multiinquilino) y versión para iOS.
  \item Frecuencia de captura adaptativa según el movimiento, para ahorrar batería.
\end{itemize}

% ════════════════════════════════════════════════════════════════════════════
\section*{Referencias}
\addcontentsline{toc}{section}{Referencias}
\begin{enumerate}[label={[\arabic*]}, leftmargin=*]
  \item Flutter. \emph{Documentación oficial}. \url{https://docs.flutter.dev}
  \item Android Developers. \emph{Foreground services} y \emph{Background location limits}.
        \url{https://developer.android.com/develop/background-work/services/fgs}
  \item Amazon Web Services. \emph{Amazon EC2} y \emph{Amazon RDS for PostgreSQL User Guide}.
        \url{https://docs.aws.amazon.com}
  \item Socket.IO. \emph{Documentation}. \url{https://socket.io/docs/v4/}
  \item Leaflet. \emph{An open-source JavaScript library for interactive maps}. \url{https://leafletjs.com}
  \item Ainsworth, B. E. et al. (2011). \emph{Compendium of Physical Activities}. Medicine \& Science in
        Sports \& Exercise, 43(8), 1575--1581.
  \item Kalman, R. E. (1960). A New Approach to Linear Filtering and Prediction Problems.
        \emph{Journal of Basic Engineering}, 82(1), 35--45.
  \item Sinnott, R. W. (1984). Virtues of the Haversine. \emph{Sky and Telescope}, 68(2), 159.
\end{enumerate}

\end{document}
